\section{Proposed formative user-study protocol}

This appendix develops the study proposed in the April interim report. It is a prospective protocol, not a completed study or an ethics approval. Recruitment should begin only after the institution's required review and after exercise content is suitable for the tasks. A minimum of five adult participants may support formative usability observations, but cannot provide representative clinical-effectiveness evidence. Participation should not require disclosure of a traumatic history or a diagnosis; any self-identification criteria require clear consent and data minimisation.

Participants would receive a plain-language explanation of purpose, activities, voluntary withdrawal, data handling, and the application's limits. Tasks would use synthetic account data and optional non-sensitive check-in choices. The facilitator would explain how to pause or stop, and participants could skip any practice. No recovery key from a participant's real journal should be collected. A demonstration vault or synthetic entry can be used to assess comprehension of setup and recovery without exposing personal reflections.

The task sequence would cover onboarding, completing a check-in, selecting a comfort preference, reading a recommendation reason, browsing an alternative practice, opening a learning resource, saving a synthetic reflection, locking the journal, and finding history or account controls. Observations would record completion, errors, requests for help, navigation confusion, and understanding of restrictions. Relevance ratings would be collected separately from comfort and perceived usefulness, so a relevant-looking suggestion is not automatically coded as beneficial.

A comparator, if included, should use the same eligible and reviewed content pool and differ only in selection method. Counterbalanced order can reduce simple learning effects, but the small study still remains formative. The study should not reproduce the offline random baseline's missing suitability rules in participant-facing tasks. Any accidental restriction violation should be recorded and investigated rather than hidden in an average relevance score.

A standard usability instrument can supplement task observations, provided its wording and scoring are reported accurately. Interviews should ask what participants expected, which explanations were useful, where control felt insufficient, and what made journal recovery confusing. Analysis can combine task counts with descriptive scores and transparent thematic summaries. In a small sample, individual difficulties are informative; a single aggregate satisfaction score should not conceal them. No clinical symptom-improvement claim follows from a short task session.

The final study report should state recruitment method, participant count, task protocol, consent and withdrawal handling, instrument version, missing responses, and analytic approach. Findings should distinguish observed behaviour from participant opinion and researcher interpretation. Retention should follow the approved research arrangement. Completing this protocol would address the original study objective and provide evidence for interface iteration, while stronger outcome claims would still require a separate research design.
